\section{The proposed hybrid algorithm}
\label{sec:algorithm}

We first state a proof-oriented algorithm.  Its Phase II uses $\omega=1$;
optimal over-relaxation and Chebyshev are treated later as optional safeguarded
or experimental variants.

\subsection{Objective-gap handoff}

Define the target handoff gap
\begin{equation}
  \Delta_{\rm sw}
  :=\frac{\alpha^{3/2}\eps}{1+\alpha}.
  \label{eq:Delta-switch}
\end{equation}
Run AESP until the deterministic upper bound in
\cref{eq:aesp-outer-bound} satisfies
\begin{equation}
  U_J\leq\Delta_{\rm sw}.
  \label{eq:objective-switch}
\end{equation}
This rule is theoretical but computable: it does not require $f^\star$.
After the switch, discard all Catalyst momentum and run local Gauss--Seidel
SOR at the final coordinate threshold \cref{eq:final-active}.

\begin{algorithm}[t]
\caption{Proof-oriented AESP--LOCSOR hybrid}
\label{alg:hybrid}
\begin{algorithmic}[1]
\REQUIRE $G,s,\alpha\in(0,1/2),\eps>0$; AESP inner method
         $\mathcal M\in\{\LOCGD,\LOCAPPR\}$
\STATE Set $x^{(0)}=y^{(0)}=0$, $\eta=1-2\alpha$.
\STATE Define $U_t$ by \cref{eq:aesp-outer-bound} and
       $\Delta_{\rm sw}$ by \cref{eq:Delta-switch}.
\FOR{$t=1,2,\ldots$}
  \STATE Use $\mathcal M$ to return $x^{(t)}$ satisfying the AESP
         inner objective requirement for stage $t$.
  \IF{$\norm{D^{-1/2}g(x^{(t)})}_\infty\leq\alpha\eps$}
    \RETURN $\widehat\pi=D^{1/2}x^{(t)}$.
  \ENDIF
  \IF{$U_t\leq\Delta_{\rm sw}$}
    \STATE Set $x\leftarrow x^{(t)}$ and discard the momentum state.
    \STATE \textbf{break}.
  \ENDIF
  \STATE $y^{(t)}\leftarrow x^{(t)}+
    \beta_{\rm A}(x^{(t)}-x^{(t-1)})$.
\ENDFOR
\WHILE{there exists $u$ with $\abs{g_u(x)}\geq\alpha\eps\sqrt{d_u}$}
  \STATE Choose any active coordinate $u$.
  \STATE $x\leftarrow x-\dfrac{2g_u(x)}{1+\alpha}e_u$.
\ENDWHILE
\RETURN $\widehat\pi=D^{1/2}x$.
\end{algorithmic}
\end{algorithm}

\begin{remark}[A coarser AESP target]
The switch target is much weaker than solving the final PageRank problem by
AESP alone.  If AESP is parameterized through a coordinate-error target
$\eps_{\rm burn}$ satisfying
$\alpha\eps_{\rm burn}^2/2=\Delta_{\rm sw}$, then
\begin{equation}
  \eps_{\rm burn}
  =\sqrt{\frac{2\sqrt\alpha\eps}{1+\alpha}}.
  \label{eq:burn-eps}
\end{equation}
For the high-accuracy regime $\eps\ll\sqrt\alpha$, this is substantially
larger than the final $\eps$.
\end{remark}

\subsection{Alternative residual-mass handoff}

Define the original weighted gradient mass
\begin{equation}
  M(x):=\norm{D^{1/2}g(x)}_1=\norm{q(x)}_1.
  \label{eq:mass-def}
\end{equation}
A second switch rule is
\begin{equation}
  M(x^{(J)})\leq\theta\alpha^{3/2},
  \qquad \theta=\bigO(1).
  \label{eq:mass-switch}
\end{equation}
This rule gives a stronger tail-locality theorem in
\cref{thm:mass-tail}, but the current AESP theory does not guarantee that the
threshold is reached with graph-independent local work.  It is therefore best
viewed as an online diagnostic and an alternative theorem hypothesis.

\section{Unconditional convergence and tail bounds}
\label{sec:tail}

\subsection{Exact coordinate descent identity}

Since $G$ is simple, $Q_{uu}=(1+\alpha)/2$.  A SOR coordinate step with
relaxation $\omega$ is
\begin{equation}
  x^+=x-\frac{\omega g_u(x)}{Q_{uu}}e_u
      =x-\frac{2\omega}{1+\alpha}g_u(x)e_u.
  \label{eq:sor-coordinate}
\end{equation}

\begin{lemma}[Exact objective decrease]
\label{lem:sor-decrease}
For every $x$, coordinate $u$, and $\omega\in\R$,
\begin{equation}
  f(x)-f(x^+)
  =\frac{\omega(2-\omega)}{1+\alpha}g_u(x)^2.
  \label{eq:sor-exact-decrease}
\end{equation}
In particular, the objective decreases strictly whenever
$g_u(x)\neq0$ and $0<\omega<2$.
\end{lemma}

\begin{proof}
For $t=-\omega g_u(x)/Q_{uu}$, the quadratic expansion gives
\[
  f(x+te_u)
  =f(x)+t g_u(x)+\frac12Q_{uu}t^2.
\]
Substituting $t$ yields
\[
  f(x)-f(x^+)
  =\frac{\omega(2-\omega)}{2Q_{uu}}g_u(x)^2
  =\frac{\omega(2-\omega)}{1+\alpha}g_u(x)^2.
\]
\end{proof}

\begin{theorem}[Finite convergence from every handoff]
\label{thm:finite-tail}
Let $x_0\in\R^n$ be arbitrary, including a signed AESP handoff point.  Fix
$\omega\in(0,2)$.  Repeatedly apply \cref{eq:sor-coordinate} to any coordinate
satisfying \cref{eq:final-active}.  Then the process terminates after finite
degree-weighted work, and its output $\widehat\pi=D^{1/2}x$ satisfies
\[
  \norm{D^{-1}(\widehat\pi-\pi)}_\infty\leq\eps.
\]
Moreover,
\begin{equation}
  W_{\rm tail}(x_0;\omega)
  \leq
  \frac{(1+\alpha)(f(x_0)-f^\star)}
       {\omega(2-\omega)\alpha^2\eps^2}.
  \label{eq:generic-tail-work}
\end{equation}
\end{theorem}

\begin{proof}
An active coordinate satisfies
$g_u(x)^2\geq\alpha^2\eps^2d_u$.  By
\cref{lem:sor-decrease}, its update decreases $f$ by at least
\[
  \frac{\omega(2-\omega)\alpha^2\eps^2}{1+\alpha}d_u.
\]
Summing over all updates and using $f(x)\geq f^\star$ gives
\cref{eq:generic-tail-work}.  Hence the total charged degree is finite, so no
active coordinate can be processed indefinitely.  At termination the test in
\cref{eq:certificate-stop} holds, and \cref{lem:certificate} gives the stated
PageRank error.
\end{proof}

\begin{corollary}[Unconditional convergence of the hybrid]
\label{cor:hybrid-converges}
Suppose the AESP phase executes any finite number of valid outer stages and
then hands off to local SOR with fixed $\omega\in(0,2)$.  The resulting hybrid
terminates and satisfies the final PageRank certificate, independently of any
active-set locality assumption.
\end{corollary}

This answers the basic correctness question affirmatively.  What remains is
the degree-weighted work of the AESP burn-in and the SOR tail.

\subsection{Objective-gap handoff gives an accelerated-size tail}

\begin{theorem}[Objective-gap tail theorem]
\label{thm:objective-tail}
Run local SOR with $\omega=1$ from a handoff point $x_J$.  If
\begin{equation}
  f(x_J)-f^\star
  \leq\frac{\alpha^{3/2}\eps}{1+\alpha},
  \label{eq:objective-tail-assumption}
\end{equation}
then
\begin{equation}
  W_{\rm tail}\leq\frac{1}{\sqrt\alpha\eps}.
  \label{eq:objective-tail-result}
\end{equation}
\end{theorem}

\begin{proof}
Set $\omega=1$ in \cref{eq:generic-tail-work} and substitute
\cref{eq:objective-tail-assumption}:
\[
  W_{\rm tail}
  \leq
  \frac{(1+\alpha)(f(x_J)-f^\star)}{\alpha^2\eps^2}
  \leq\frac{1}{\sqrt\alpha\eps}.
\]
\end{proof}

\begin{remark}
The tail is not spectrally accelerated in this theorem.  The factor
$1/\sqrt\alpha$ is supplied by the quality of the AESP handoff.  Safe local
coordinate descent merely converts the remaining objective budget into the
final $1/\eps$ accuracy dependence.
\end{remark}

\subsection{Signed weighted-$\ell_1$ tail theorem}

The objective-gap argument gives total tail work but does not itself bound
each active-set volume.  The weighted gradient mass does.

For $q=D^{1/2}g$ and a SOR step at $u$, define
$\chi=(1-\alpha)/(1+\alpha)$.  Multiplying the gradient recurrence by
$D^{1/2}$ gives
\begin{equation}
  q^+=q-\omega q_u e_u+\chi\omega q_uAD^{-1}e_u.
  \label{eq:q-sor-update}
\end{equation}
For $\omega=1$, this becomes
\begin{equation}
  q_u^+=0,
  \qquad
  q_v^+=q_v+\chi\frac{q_u}{d_u}
  \quad(v\sim u).
  \label{eq:q-gs-update}
\end{equation}

\begin{lemma}[Signed weighted-$\ell_1$ decrease at $\omega=1$]
\label{lem:signed-l1}
For every signed vector $q$ and every coordinate update
\cref{eq:q-gs-update},
\begin{equation}
  \norm{q^+}_1
  \leq
  \norm{q}_1-\frac{2\alpha}{1+\alpha}\abs{q_u}.
  \label{eq:signed-l1-drop}
\end{equation}
\end{lemma}

\begin{proof}
The column $AD^{-1}e_u$ is nonnegative and has $\ell_1$ norm one.  Therefore
\[
  \norm{q^+}_1
  \leq
  \norm{q}_1-\abs{q_u}
  +\chi\abs{q_u}
  =\norm{q}_1-(1-\chi)\abs{q_u},
\]
and $1-\chi=2\alpha/(1+\alpha)$.
\end{proof}

\begin{theorem}[Residual-mass tail theorem]
\label{thm:mass-tail}
Let $x_0$ be any handoff point and run local SOR with $\omega=1$.  Then
\begin{equation}
  W_{\rm tail}
  \leq
  \frac{1+\alpha}{2\alpha^2\eps}M(x_0).
  \label{eq:mass-tail-general}
\end{equation}
If $M(x_0)\leq\theta\alpha^{3/2}$, then
\begin{equation}
  W_{\rm tail}
  \leq
  \frac{\theta(1+\alpha)}{2\sqrt\alpha\eps},
  \label{eq:mass-tail-accelerated}
\end{equation}
and every active set $S_k$ in the tail satisfies
\begin{equation}
  \vol(S_k)
  \leq\frac{\theta\sqrt\alpha}{\eps}.
  \label{eq:mass-tail-volume}
\end{equation}
\end{theorem}

\begin{proof}
An active update has $\abs{q_u}\geq\alpha\eps d_u$.  Combining this with
\cref{eq:signed-l1-drop} gives
\[
  M(x)-M(x^+)
  \geq\frac{2\alpha^2\eps}{1+\alpha}d_u.
\]
Summing and using nonnegativity of $M$ yields
\cref{eq:mass-tail-general}; the handoff assumption gives
\cref{eq:mass-tail-accelerated}.  Moreover, at any tail iterate,
\[
  \alpha\eps\vol(S_k)
  \leq\sum_{u\in S_k}\abs{q_u}
  \leq M(x_k)
  \leq M(x_0),
\]
which proves \cref{eq:mass-tail-volume}.
\end{proof}
